\documentclass[10pt,a4paper]{amsart}
\usepackage{amsmath,amssymb,amsthm,graphicx}
% For commutative diagrams we use Xy-pic
\usepackage[all]{xy}

% We use 2cell for 2-commutative diagrams.
\xyoption{2cell}
% Theorem environments.
%
\theoremstyle{plain}
\newtheorem{theorem}[subsection]{Theorem}
\newtheorem{proposition}[subsection]{Proposition}
\newtheorem{lemma}[subsection]{Lemma}

\theoremstyle{definition}
\newtheorem{definition}[subsection]{Definition}
\newtheorem{example}[subsection]{Example}
\newtheorem{exercise}[subsection]{Exercise}
\newtheorem{situation}[subsection]{Situation}

\theoremstyle{remark}
\newtheorem{remark}[subsection]{Remark}
\newtheorem{remarks}[subsection]{Remarks}

\numberwithin{equation}{subsection}

% Macros
%
\def\lim{\mathop{\mathrm{lim}}\nolimits}
\def\colim{\mathop{\mathrm{colim}}\nolimits}
\def\Spec{\mathop{\mathrm{Spec}}}
\def\Hom{\mathop{\mathrm{Hom}}\nolimits}
\def\Ext{\mathop{\mathrm{Ext}}\nolimits}
\def\SheafHom{\mathop{\mathcal{H}\!\mathit{om}}\nolimits}
\def\SheafExt{\mathop{\mathcal{E}\!\mathit{xt}}\nolimits}
\def\Sch{\mathit{Sch}}
\def\Mor{\mathop{\mathrm{Mor}}\nolimits}
\def\Ob{\mathop{\mathrm{Ob}}\nolimits}
\def\Sh{\mathop{\mathit{Sh}}\nolimits}
\def\NL{\mathop{N\!L}\nolimits}
\def\CH{\mathop{\mathrm{CH}}\nolimits}
\def\proetale{{pro\text{-}\acute{e}tale}}
\def\etale{{\acute{e}tale}}
\def\QCoh{\mathit{QCoh}}
\def\Ker{\mathop{\mathrm{Ker}}}
\def\Im{\mathop{\mathrm{Im}}}
\def\Coker{\mathop{\mathrm{Coker}}}
\def\Coim{\mathop{\mathrm{Coim}}}

% Boxtimes
%
\DeclareMathSymbol{\boxtimes}{\mathbin}{AMSa}{"02}

%
% Macros for moduli stacks/spaces
%
\def\QCohstack{\mathcal{QC}\!\mathit{oh}}
\def\Cohstack{\mathcal{C}\!\mathit{oh}}
\def\Spacesstack{\mathcal{S}\!\mathit{paces}}
\def\Quotfunctor{\mathrm{Quot}}
\def\Hilbfunctor{\mathrm{Hilb}}
\def\Curvesstack{\mathcal{C}\!\mathit{urves}}
\def\Polarizedstack{\mathcal{P}\!\mathit{olarized}}
\def\Complexesstack{\mathcal{C}\!\mathit{omplexes}}
% \Pic is the operator that assigns to X its picard group, usage \Pic(X)
% \Picardstack_{X/B} denotes the Picard stack of X over B
% \Picardfunctor_{X/B} denotes the Picard functor of X over B
\def\Pic{\mathop{\mathrm{Pic}}\nolimits}
\def\Picardstack{\mathcal{P}\!\mathit{ic}}
\def\Picardfunctor{\mathrm{Pic}}
\def\Deformationcategory{\mathcal{D}\!\mathit{ef}}
\usepackage{cleveref}
\crefname{originalequation}{original equation}{original equations}
\Crefname{originalequation}{Original equation}{Original equations}
\makeatletter
\newlabel{equation-compare-big-small}{{070U}{source}{}{sourcepointer-1}{}}
\newlabel{equation-compare-big-small@cref}{{[originalequation][0][]070U}{[1][0][]source}{}{}{}}
\makeatother
\crefname{originalreference}{original reference}{original references}
\Crefname{originalreference}{Original reference}{Original references}
\makeatletter
\newlabel{lemma-compare-injectives}{{0758}{source}{}{sourcepointer-2}{}}
\newlabel{lemma-compare-injectives@cref}{{[originalreference][0][]0758}{[1][0][]source}{}{}{}}
\makeatother
\crefname{originallemma}{original lemma}{original lemmas}
\Crefname{originallemma}{Original lemma}{Original lemmas}
\makeatletter
\newlabel{sites-cohomology-lemma-derived-pushforward-composition}{{0D6E}{source}{}{sourcepointer-3}{}}
\newlabel{sites-cohomology-lemma-derived-pushforward-composition@cref}{{[originallemma][0][]0D6E}{[1][0][]source}{}{}{}}
\makeatother
\crefname{originallemma}{original lemma}{original lemmas}
\Crefname{originallemma}{Original lemma}{Original lemmas}
\makeatletter
\newlabel{sites-cohomology-lemma-localize-cartesian-square}{{0EZ0}{source}{}{sourcepointer-4}{}}
\newlabel{sites-cohomology-lemma-localize-cartesian-square@cref}{{[originallemma][0][]0EZ0}{[1][0][]source}{}{}{}}
\makeatother
\crefname{originallemma}{original lemma}{original lemmas}
\Crefname{originallemma}{Original lemma}{Original lemmas}
\makeatletter
\newlabel{sites-cohomology-lemma-localize-cartesian-square-modules}{{0FN5}{source}{}{sourcepointer-5}{}}
\newlabel{sites-cohomology-lemma-localize-cartesian-square-modules@cref}{{[originallemma][0][]0FN5}{[1][0][]source}{}{}{}}
\makeatother
\crefname{originallemma}{original lemma}{original lemmas}
\Crefname{originallemma}{Original lemma}{Original lemmas}
\makeatletter
\newlabel{topologies-lemma-morphism-big-small-cartesian-diagram-etale}{{0DDA}{source}{}{sourcepointer-6}{}}
\newlabel{topologies-lemma-morphism-big-small-cartesian-diagram-etale@cref}{{[originallemma][0][]0DDA}{[1][0][]source}{}{}{}}
\makeatother
\title{The big pushforward restricts to the small one, and the comparison is stable under base change}
\author{The Stacks Project Authors}
\date{stacks-project:etale-cohomology.tex:local-snapshot-20261011; GNU FDL 1.2 or later}
\begin{document}
\maketitle
\noindent\textit{Editorial note.} All delivered mathematical text is the author's own, copied from the pinned Stacks Project chapter file etale-cohomology.tex; no statement or proof step was rewritten, reordered or removed.\par\medskip
\noindent\textit{Editorial note.} The counted proof is the author's complete argument: the proof is written out in full for all six parts: the K-injective computation, the two localisation statements, the composition of ringed pushforwards with exactness, and the base changed module case.\par\medskip
\noindent\textit{Editorial note.} The label lemma-compare-higher-direct-image names the counted statement, and the proof environment immediately following it in the source is delivered verbatim as its proof.\par\medskip
\noindent\textit{Editorial note.} The counted statement is carried by the author's own lemma environment in the chapter source.\par\medskip
\noindent\textit{Editorial note.} The statement and proof are copied from the pinned chapter file; the seven cited inputs are kept as external pointers with their canonical Stacks tags.\par\medskip
\noindent\textit{Allowed standard prerequisites (editorial).}\begin{itemize}
\item Injectivity of the restrictions from the big to the small etale site. Same chapter, the lemma stating that the restriction of an injective or K-injective abelian sheaf on the big site stays injective on the small one. Supplies the computation of parts (1) and (3) through K-injective representatives.
\item The cartesian diagram of topologies attached to an object of the big site. Topologies, the lemma giving the commutative diagram of the big and small etale topoi attached to a morphism of schemes. Supplies the diagram used in part (3).
\item Localisation of cartesian squares, with and without modules. Cohomology on Sites, the two lemmas identifying the localised pushforward with the pushforward of the localised complex. Supplies parts (5) and (6) directly.
\end{itemize}\medskip
\ifdefined\hypertarget\hypertarget{sourcepointer-1}{}\fi
\noindent\textit{Source reference.} Original equation 070U: sources/stacks-etale-cohomology/pinned-original-source.tex, source lines 21187--21187. Prerequisite pointer invoked at proof line(s) 21238; it names a standard fact of the same chapter that the counted argument uses, and every step of the counted proof built on it is retained verbatim in the delivered body.\par\medskip
\ifdefined\hypertarget\hypertarget{sourcepointer-2}{}\fi
\noindent\textit{Source reference.} Original reference 0758: sources/stacks-etale-cohomology/pinned-original-source.tex, source lines 21102--21102. Prerequisite pointer invoked at proof line(s) 21237; it names a standard fact of the same chapter that the counted argument uses, and every step of the counted proof built on it is retained verbatim in the delivered body.\par\medskip
\ifdefined\hypertarget\hypertarget{sourcepointer-3}{}\fi
\noindent\textit{Source reference.} Original lemma 0D6E: sources/stacks-chow/tags.txt, source lines 16519--16519. Prerequisite pointer invoked at proof line(s) 21261; it names a standard fact of a sibling chapter that the counted argument uses, and every step of the counted proof built on it is retained verbatim in the delivered body.\par\medskip
\ifdefined\hypertarget\hypertarget{sourcepointer-4}{}\fi
\noindent\textit{Source reference.} Original lemma 0EZ0: sources/stacks-chow/tags.txt, source lines 18710--18710. Prerequisite pointer invoked at proof line(s) 21249; it names a standard fact of a sibling chapter that the counted argument uses, and every step of the counted proof built on it is retained verbatim in the delivered body.\par\medskip
\ifdefined\hypertarget\hypertarget{sourcepointer-5}{}\fi
\noindent\textit{Source reference.} Original lemma 0FN5: sources/stacks-chow/tags.txt, source lines 19555--19555. Prerequisite pointer invoked at proof line(s) 21253; it names a standard fact of a sibling chapter that the counted argument uses, and every step of the counted proof built on it is retained verbatim in the delivered body.\par\medskip
\ifdefined\hypertarget\hypertarget{sourcepointer-6}{}\fi
\noindent\textit{Source reference.} Original lemma 0DDA: sources/stacks-chow/tags.txt, source lines 16760--16760. Prerequisite pointer invoked at proof line(s) 21244; it names a standard fact of a sibling chapter that the counted argument uses, and every step of the counted proof built on it is retained verbatim in the delivered body.\par\medskip
\begin{lemma}
\label{lemma-compare-higher-direct-image}
Let $f : T \to S$ be a morphism of schemes.
\begin{enumerate}
\item For $K$ in $D((\Sch/T)_\etale)$ we have
$
(Rf_{big, *}K)|_{S_\etale} = Rf_{small, *}(K|_{T_\etale})
$
in $D(S_\etale)$.
\item For $K$ in $D((\Sch/T)_\etale, \mathcal{O})$ we have
$
(Rf_{big, *}K)|_{S_\etale} = Rf_{small, *}(K|_{T_\etale})
$
in $D(\textit{Mod}(S_\etale, \mathcal{O}_S))$.
\end{enumerate}
More generally, let $g : S' \to S$ be an object of $(\Sch/S)_\etale$.
Consider the fibre product
$$
\xymatrix{
T' \ar[r]_{g'} \ar[d]_{f'} & T \ar[d]^f \\
S' \ar[r]^g & S
}
$$
Then
\begin{enumerate}
\item[(3)] For $K$ in $D((\Sch/T)_\etale)$ we have
$i_g^{-1}(Rf_{big, *}K) = Rf'_{small, *}(i_{g'}^{-1}K)$
in $D(S'_\etale)$.
\item[(4)] For $K$ in $D((\Sch/T)_\etale, \mathcal{O})$ we have
$i_g^*(Rf_{big, *}K) = Rf'_{small, *}(i_{g'}^*K)$
in $D(\textit{Mod}(S'_\etale, \mathcal{O}_{S'}))$.
\item[(5)] For $K$ in $D((\Sch/T)_\etale)$ we have
$g_{big}^{-1}(Rf_{big, *}K) = Rf'_{big, *}((g'_{big})^{-1}K)$
in $D((\Sch/S')_\etale)$.
\item[(6)] For $K$ in $D((\Sch/T)_\etale, \mathcal{O})$ we have
$g_{big}^*(Rf_{big, *}K) = Rf'_{big, *}((g'_{big})^*K)$
in $D(\textit{Mod}(S'_\etale, \mathcal{O}_{S'}))$.
\end{enumerate}
\end{lemma}

\begin{proof}
Part (1) follows from
Lemma \ref{lemma-compare-injectives}
and (\ref{equation-compare-big-small})
on choosing a K-injective complex of abelian sheaves representing $K$.

\medskip\noindent
Part (3) follows from Lemma \ref{lemma-compare-injectives}
and Topologies, Lemma
\ref{topologies-lemma-morphism-big-small-cartesian-diagram-etale}
on choosing a K-injective complex of abelian sheaves representing $K$.

\medskip\noindent
Part (5) is Cohomology on Sites, Lemma
\ref{sites-cohomology-lemma-localize-cartesian-square}.

\medskip\noindent
Part (6) is Cohomology on Sites, Lemma
\ref{sites-cohomology-lemma-localize-cartesian-square-modules}.

\medskip\noindent
Part (2) can be proved as follows. Above we have seen
that $\pi_S \circ f_{big} = f_{small} \circ \pi_T$ as morphisms
of ringed sites. Hence we obtain
$R\pi_{S, *} \circ Rf_{big, *} = Rf_{small, *} \circ R\pi_{T, *}$
by Cohomology on Sites, Lemma
\ref{sites-cohomology-lemma-derived-pushforward-composition}.
Since the restriction functors $\pi_{S, *}$ and $\pi_{T, *}$
are exact, we conclude.

\medskip\noindent
Part (4) follows from part (6) and part (2) applied to $f' : T' \to S'$.
\end{proof}

\end{document}
